\section{Introduction}
\label{sec:introduction}

Recommendation from implicit feedback differs from rating prediction in a way
that is easy to state and easy to get wrong. We never observe a negative
signal. A title a user has not watched may be a title they would dislike, or
simply one they have not encountered. Optimizing a squared error against a
zero-filled matrix therefore encodes an assumption that is false by
construction: that absence of interaction means dislike.

Bayesian Personalized Ranking~\cite{rendle2009bpr} replaces that assumption
with a weaker and more defensible one. Rather than predicting an absolute
score, it assumes only that an observed item is preferred to an unobserved
one, and optimizes the probability of that \emph{pairwise} ordering. The
learning signal is a triplet $(u, i, j)$ where $i$ was consumed and $j$ was
not, and the objective rewards $\hat{x}_{ui} > \hat{x}_{uj}$.

\subsection{Scope and claims}

This report accompanies a working repository rather than a novel method. The
claims we make are:

\begin{enumerate}
  \item The pipeline is leakage-free. Every aggregate feature is computed
        point-in-time, and this is enforced by tests that inject a future
        event and assert no feature changes (Section~\ref{sec:features}).
  \item The reported numbers are honest. Hyperparameters for \emph{every}
        model, baselines included, were selected on a validation partition;
        the test partition was read once (Section~\ref{sec:evaluation}).
  \item The served ranking equals the evaluated ranking. Training/serving
        skew is asserted away rather than assumed away
        (Section~\ref{sec:evaluation}).
\end{enumerate}

We explicitly do \emph{not} claim that BPR is the best model for this dataset.
Section~\ref{sec:results} shows it is not.

\subsection{Why the negative result is kept}

Dacrema et al.~\cite{dacrema2019} reproduced eighteen neural recommendation
papers and found that most were beaten by well-tuned non-neural baselines,
largely because the baselines in the original papers had not been tuned at
all. The cheapest way to produce a flattering result for a flagship model is
to compare it against an untuned competitor. We tune the baselines on the same
grid budget as the flagship model, and report the outcome as measured.

\subsection{Repository}

The implementation, infrastructure, and the scripts that generate every table
in this document are available in the accompanying repository. The results
table (Table~\ref{tab:results}) is emitted by
\texttt{scripts/export\_results\_table.py} and the feature catalogue
(Table~\ref{tab:features}) by \texttt{scripts/export\_feature\_tables.py},
so the document cannot drift away from the code.
